% Options for packages loaded elsewhere
\PassOptionsToPackage{unicode}{hyperref}
\PassOptionsToPackage{hyphens}{url}
\PassOptionsToPackage{dvipsnames,svgnames,x11names}{xcolor}
%
\documentclass[
  10.5pt,
]{article}
\usepackage{amsmath,amssymb}
\usepackage{iftex}
\ifPDFTeX
  \usepackage[T1]{fontenc}
  \usepackage[utf8]{inputenc}
  \usepackage{textcomp} % provide euro and other symbols
\else % if luatex or xetex
  \usepackage{unicode-math} % this also loads fontspec
  \defaultfontfeatures{Scale=MatchLowercase}
  \defaultfontfeatures[\rmfamily]{Ligatures=TeX,Scale=1}
\fi
\usepackage{lmodern}
\ifPDFTeX\else
  % xetex/luatex font selection
\fi
% Use upquote if available, for straight quotes in verbatim environments
\IfFileExists{upquote.sty}{\usepackage{upquote}}{}
\IfFileExists{microtype.sty}{% use microtype if available
  \usepackage[]{microtype}
  \UseMicrotypeSet[protrusion]{basicmath} % disable protrusion for tt fonts
}{}
\makeatletter
\@ifundefined{KOMAClassName}{% if non-KOMA class
  \IfFileExists{parskip.sty}{%
    \usepackage{parskip}
  }{% else
    \setlength{\parindent}{0pt}
    \setlength{\parskip}{6pt plus 2pt minus 1pt}}
}{% if KOMA class
  \KOMAoptions{parskip=half}}
\makeatother
\usepackage{xcolor}
\usepackage[margin=0.85in]{geometry}
\usepackage{listings}
\newcommand{\passthrough}[1]{#1}
\lstset{defaultdialect=[5.3]Lua}
\lstset{defaultdialect=[x86masm]Assembler}

\lstset{
  basicstyle=\ttfamily\small,
  breaklines=true,
  breakatwhitespace=false,
  columns=fullflexible,
  keepspaces=true,
  showstringspaces=false,
  frame=single,
  framerule=0.2pt,
  xleftmargin=0.8em,
  xrightmargin=0.8em,
  aboveskip=0.8em,
  belowskip=0.8em
}
\setlength{\emergencystretch}{3em} % prevent overfull lines
\providecommand{\tightlist}{%
  \setlength{\itemsep}{0pt}\setlength{\parskip}{0pt}}
\setcounter{secnumdepth}{5}
\ifLuaTeX
  \usepackage{selnolig}  % disable illegal ligatures
\fi
\usepackage{bookmark}
\IfFileExists{xurl.sty}{\usepackage{xurl}}{} % add URL line breaks if available
\urlstyle{same}
\hypersetup{
  colorlinks=true,
  linkcolor={Maroon},
  filecolor={Maroon},
  citecolor={Blue},
  urlcolor={Blue},
  pdfcreator={LaTeX via pandoc}}

\title{MiNES Variance-Weighted PMF Fusion: Implementation Specification}
\author{Adaptive Nonequilibrium Sampling Notes}
\date{May 5, 2026}

\begin{document}
\maketitle

{
\hypersetup{linkcolor=}
\setcounter{tocdepth}{3}
\tableofcontents
}
\section{Task: Implement a clean MiNES-only workflow with
variance-weighted PMF
fusion}\label{task-implement-a-clean-mines-only-workflow-with-variance-weighted-pmf-fusion}

You are working in the current repository for adaptive nonequilibrium
sampling. The existing \passthrough{\lstinline!adaptive\_methods.py!} is
too broad and does not implement the intended MiNES algorithm cleanly.
Do \textbf{not} modify \passthrough{\lstinline!adaptive\_methods.py!}
unless absolutely necessary. Instead, create a new MiNES-only
implementation.

\subsection{Deliverables}\label{deliverables}

Create three new files:

\begin{enumerate}
\def\labelenumi{\arabic{enumi}.}
\tightlist
\item
  \passthrough{\lstinline!scripts/mines\_variance\_fusion.py!}
\item
  \passthrough{\lstinline!scripts/run\_mines\_variance\_fusion.sh!}
\item
  \passthrough{\lstinline!notebooks/mines\_variance\_fusion\_visualization.ipynb!}
\end{enumerate}

The new Python file should only handle MiNES. It should not include AUS
logic.

\begin{center}\rule{0.5\linewidth}{0.5pt}\end{center}

\section{\texorpdfstring{1. New Python file:
\texttt{scripts/mines\_variance\_fusion.py}}{1. New Python file: scripts/mines\_variance\_fusion.py}}\label{new-python-file-scriptsmines_variance_fusion.py}

\subsection{Goal}\label{goal}

Implement MiNES as a chain of EQ clusters and NEQ segments, where every
local reconstruction produces:

\begin{lstlisting}[language=Python]
PMF patch:
    pmf(x)
    variance(x)
    coverage_mask(x)
    unknown additive offset
\end{lstlisting}

The final PMF should be obtained by \textbf{global inverse-variance
weighted fusion of all local PMF patches}, not by hard stitching.

The high-level workflow is:

\begin{enumerate}
\def\labelenumi{\arabic{enumi}.}
\tightlist
\item
  Perform EQ simulations at endpoint ensembles
  \passthrough{\lstinline!L0!} and \passthrough{\lstinline!R0!}.
\item
  Perform NEQ switching between \passthrough{\lstinline!L0!} and
  \passthrough{\lstinline!R0!}.
\item
  Use the selected leaping method to decide child ensembles
  \passthrough{\lstinline!L1!}, \passthrough{\lstinline!R1!}, then
  perform EQ simulations.
\item
  For every neighboring pair of ensembles/clusters:

  \begin{itemize}
  \tightlist
  \item
    If JS divergence is below the threshold, merge them into an EQ
    cluster and reconstruct an EQ MBAR patch.
  \item
    If JS divergence is above the threshold, construct an NEQ MTS patch
    between them.
  \end{itemize}
\item
  Continue growing children until the left and right frontiers cross or
  meet.
\item
  After crossing, enter the rescue phase.
\item
  In rescue, target the grid bin with the largest global fused variance.
\item
  Add a rescue EQ window or NEQ bridge depending on where the
  maximum-variance bin lies.
\item
  Reconstruct all patches and refit the global PMF after every rescue
  round.
\end{enumerate}

\begin{center}\rule{0.5\linewidth}{0.5pt}\end{center}

\section{Required command-line
interface}\label{required-command-line-interface}

The script should be executable as:

\begin{lstlisting}[language=bash]
python scripts/mines_variance_fusion.py \
  --system-root <path> \
  --bin <path-to-neq_sim> \
  --seed 123 \
  --label mines_variance_fusion \
  --total-budget-steps 200000 \
  --t-neq 5000 \
  --n-neq-traj 100
\end{lstlisting}

Arguments:

\begin{lstlisting}[language=Python]
--system-root              required
--bin                      required, path to compiled simulation binary
--seed                     required int
--label                    default: mines_variance_fusion
--total-budget-steps       default: 200000
--t-neq                    default: 5000
--n-neq-traj               default: 100
--n-eq-steps               default: 10000
--eq-save-every            default: 10
--tail-fraction            default: 0.9
--q-next                   default: 0.9
--alpha                    default: 2.0
--x-leap                   default: 1.5
--k-min                    default: 1.0
--k-max                    default: 50.0
--k-rescue                 default: 10.0
--js-threshold             default: 0.8
--bin-width                default: 0.1
--max-generations          default: 10
--max-rescue-rounds        default: 5
--n-bootstrap-eq           default: 64
--n-bootstrap-neq          default: 64
--variance-floor           default: 1e-6
\end{lstlisting}

Output root:

\begin{lstlisting}
<system-root>/MINES/<label>/raw/seed_<seed>/
\end{lstlisting}

\begin{center}\rule{0.5\linewidth}{0.5pt}\end{center}

\section{Core data structures}\label{core-data-structures}

Implement explicit data classes.

\subsection{\texorpdfstring{\texttt{EnsembleWindow}}{EnsembleWindow}}\label{ensemblewindow}

\begin{lstlisting}[language=Python]
@dataclass
class EnsembleWindow:
    name: str
    center_x: float
    k: float
    root: Path
    eq_file: Path
    tail_file: Path
    eq_rows: list[dict[str, str]]
    tail_rows: list[dict[str, str]]
    x_most: float
    generation: int
    side: str  # "left", "right", "rescue", "bridge"
\end{lstlisting}

\subsection{\texorpdfstring{\texttt{EQCluster}}{EQCluster}}\label{eqcluster}

\begin{lstlisting}[language=Python]
@dataclass
class EQCluster:
    name: str
    windows: list[EnsembleWindow]
    left_x: float
    right_x: float
\end{lstlisting}

\subsection{\texorpdfstring{\texttt{NEQSegment}}{NEQSegment}}\label{neqsegment}

\begin{lstlisting}[language=Python]
@dataclass
class NEQSegment:
    name: str
    left: EQCluster | EnsembleWindow
    right: EQCluster | EnsembleWindow
    root: Path
    forward_trajectories: list[list[dict[str, str]]]
    reverse_trajectories: list[list[dict[str, str]]]
    forward_path_file: Path
    reverse_path_file: Path
\end{lstlisting}

\subsection{\texorpdfstring{\texttt{PMFPatch}}{PMFPatch}}\label{pmfpatch}

\begin{lstlisting}[language=Python]
@dataclass
class PMFPatch:
    name: str
    kind: str  # "EQ_MBAR" or "NEQ_MTS"
    grid: np.ndarray
    pmf: np.ndarray
    variance: np.ndarray
    coverage_mask: np.ndarray
    source_names: list[str]
    metadata: dict
\end{lstlisting}

The final global PMF should be constructed only from
\passthrough{\lstinline!PMFPatch!} objects.

\begin{center}\rule{0.5\linewidth}{0.5pt}\end{center}

\section{Simulation helpers}\label{simulation-helpers}

Reuse the existing simulator interface style from
\passthrough{\lstinline!adaptive\_methods.py!}.

Implement:

\begin{lstlisting}[language=Python]
run_eq_window(...)
run_neq_protocol(...)
write_protocol_path(...)
read_csv_rows(...)
write_csv(...)
write_json(...)
load_json(...)
\end{lstlisting}

Use \passthrough{\lstinline!run\_context.json!} from
\passthrough{\lstinline!system\_root!}.

The simulation binary should be called with the same low-level arguments
used in the existing code:

\begin{lstlisting}
-pot
-one-dimension
-thermal_kT
-dt
-gamma
-k0
-x0
-k1
-x1
-E1
-k
-center_xy
-eq_out
-T_eq
-eq_nout
-eq_seed
-out_dir
-log
\end{lstlisting}

For NEQ, use protocol path files and single-start trajectories, similar
to the current implementation.

\begin{center}\rule{0.5\linewidth}{0.5pt}\end{center}

\section{PMF reconstruction logic}\label{pmf-reconstruction-logic}

\subsection{EQ cluster patch}\label{eq-cluster-patch}

Implement:

\begin{lstlisting}[language=Python]
build_eq_cluster_patch(cluster: EQCluster, grid: np.ndarray, ctx: dict, n_boot: int) -> PMFPatch
\end{lstlisting}

For a cluster with windows \(m = 1,\dots,M\):

\begin{enumerate}
\def\labelenumi{\arabic{enumi}.}
\tightlist
\item
  Reconstruct the base PMF by MBAR using all EQ tail samples from all
  windows in the cluster.
\item
  Bootstrap the EQ tail samples window-wise:

  \begin{itemize}
  \tightlist
  \item
    For each bootstrap replicate, resample each window's tail samples
    with replacement.
  \item
    Reconstruct the MBAR PMF.
  \end{itemize}
\item
  For each possible anchor window \(m\):

  \begin{itemize}
  \tightlist
  \item
    Align every bootstrap PMF to zero at \(x^{most}_m\).
  \item
    Compute: \[
    \mathrm{var}_m(x) = \mathrm{Var}_b[F_b(x)-F_b(x^{most}_m)]
    \]
  \end{itemize}
\item
  Define: \[
  \mathrm{var}_{EQ}(x)=\min_m \mathrm{var}_m(x)
  \]
\item
  The base PMF can be stored in any consistent gauge, e.g.~shifted to
  minimum zero. The global fusion step will determine patch offsets, so
  the local gauge does not need to be final.
\item
  The coverage mask should be finite PMF and finite variance.
\end{enumerate}

Write patch files:

\begin{lstlisting}
patches/<patch_name>_pmf.csv
patches/<patch_name>_bootstrap_variance.csv
patches/<patch_name>_summary.json
\end{lstlisting}

CSV columns:

\begin{lstlisting}
x, pmf, variance, coverage, analytic_pmf
\end{lstlisting}

Also write per-anchor variance columns if useful:

\begin{lstlisting}
var_anchor_<window_name>
\end{lstlisting}

\subsection{NEQ MTS patch}\label{neq-mts-patch}

Implement:

\begin{lstlisting}[language=Python]
build_neq_mts_patch(segment: NEQSegment, grid: np.ndarray, ctx: dict, n_boot: int) -> PMFPatch
\end{lstlisting}

For each non-overlapping neighboring pair:

\begin{enumerate}
\def\labelenumi{\arabic{enumi}.}
\tightlist
\item
  Run bidirectional NEQ.
\item
  Reconstruct PMF using bidirectional MTS.
\item
  Bootstrap the bidirectional MTS estimator.
\item
  Align each bootstrap PMF to the left anchor \(x^{most}_{left}\) and
  compute: \[
  \mathrm{var}_f(x)
  \]
\item
  Align each bootstrap PMF to the right anchor \(x^{most}_{right}\) and
  compute: \[
  \mathrm{var}_r(x)
  \]
\item
  Define: \[
  \mathrm{var}_{NEQ}(x)=\min(\mathrm{var}_f(x), \mathrm{var}_r(x))
  \]
\item
  Coverage mask should be finite PMF and finite variance.
\end{enumerate}

Write patch files:

\begin{lstlisting}
patches/<patch_name>_pmf.csv
patches/<patch_name>_bootstrap_variance.csv
patches/<patch_name>_summary.json
\end{lstlisting}

CSV columns:

\begin{lstlisting}
x, pmf, variance, var_left_anchor, var_right_anchor, coverage, analytic_pmf
\end{lstlisting}

Use existing functions from
\passthrough{\lstinline!analysis/bidirectional\_mts\_pmf.py!} where
possible:

\begin{lstlisting}[language=Python]
estimate_intermediate_reduced_free_energies
build_bidirectional_mts_pmf
bootstrap_bidirectional_mts_pmf
align_pmf_to_reference
align_pmf_to_endpoint_average_zero
\end{lstlisting}

Do not duplicate these algorithms unless necessary.

\begin{center}\rule{0.5\linewidth}{0.5pt}\end{center}

\section{Global PMF fusion}\label{global-pmf-fusion}

Implement:

\begin{lstlisting}[language=Python]
fit_global_pmf_from_patches(
    patches: list[PMFPatch],
    grid: np.ndarray,
    variance_floor: float,
    reference_x: float | None = None,
) -> tuple[np.ndarray, np.ndarray, dict]
\end{lstlisting}

The final PMF should solve:

\[
\min_{G,c}
\sum_p \sum_{x\in S_p}
\frac{(G(x)-F_p(x)-c_p)^2}{\sigma_p^2(x)+\epsilon}
\]

where:

\begin{itemize}
\tightlist
\item
  \(G(x)\) is the global PMF.
\item
  \(F_p(x)\) is patch \(p\)'s PMF.
\item
  \(c_p\) is the additive offset of patch \(p\).
\item
  \(\sigma_p^2(x)\) is the patch variance.
\item
  \(\epsilon\) is \passthrough{\lstinline!variance\_floor!}.
\end{itemize}

Fix the gauge by setting either:

\begin{lstlisting}
G(reference_x) = 0
\end{lstlisting}

or, if \passthrough{\lstinline!reference\_x is None!}:

\begin{lstlisting}
min finite G(x) = 0
\end{lstlisting}

Also estimate approximate fused variance:

\[
\mathrm{var}_{global}(x)
=
\left[
\sum_{p\in \mathcal{P}(x)}
\frac{1}{\mathrm{var}_p(x)+\epsilon}
\right]^{-1}
\]

This variance is approximate because patch errors are correlated, but it
is useful for rescue targeting.

Write:

\begin{lstlisting}
global_pmf.csv
\end{lstlisting}

with columns:

\begin{lstlisting}
x, global_pmf, global_variance, global_std, n_covering_patches, analytic_pmf
\end{lstlisting}

Also write:

\begin{lstlisting}
global_fit_summary.json
\end{lstlisting}

containing patch offsets, number of patches, RMS residuals, and gauge
reference.

\begin{center}\rule{0.5\linewidth}{0.5pt}\end{center}

\section{Adaptive MiNES growth logic}\label{adaptive-mines-growth-logic}

Implement a clean chain-based algorithm.

Initialize:

\begin{lstlisting}
left_windows = [L0]
right_windows = [R0]
clusters = [cluster_L0, cluster_R0]
segments = []
patches = []
\end{lstlisting}

Step 1:

\begin{itemize}
\tightlist
\item
  EQ sample \passthrough{\lstinline!L0!}, \passthrough{\lstinline!R0!}.
\item
  NEQ sample between \passthrough{\lstinline!L0!},
  \passthrough{\lstinline!R0!}.
\item
  Build an initial NEQ MTS patch.
\end{itemize}

Then for each generation:

\begin{enumerate}
\def\labelenumi{\arabic{enumi}.}
\tightlist
\item
  Use the chosen leaping rule to propose
  \passthrough{\lstinline!L\_next!} from the current left frontier and
  \passthrough{\lstinline!R\_next!} from the current right frontier.
\item
  Run EQ for \passthrough{\lstinline!L\_next!},
  \passthrough{\lstinline!R\_next!}.
\item
  Insert them into the chain.
\item
  Rebuild neighboring relationships.
\item
  For each neighboring pair:

  \begin{itemize}
  \tightlist
  \item
    Compute JS divergence from EQ tail samples.
  \item
    If \passthrough{\lstinline!JS < js\_threshold!}, merge into an EQ
    cluster and build one EQ MBAR patch.
  \item
    Else, run NEQ if it does not already exist and build an NEQ MTS
    patch.
  \end{itemize}
\item
  Fit global PMF from all patches.
\item
  Stop growth when the left and right frontiers cross or meet:

  \begin{itemize}
  \tightlist
  \item
    e.g.~\passthrough{\lstinline!L\_next.center\_x >= R\_next.center\_x!},
    or
  \item
    \passthrough{\lstinline!L\_next.x\_most >= R\_next.x\_most!}, or
  \item
    target-cross condition based on the chosen leaping rule.
  \end{itemize}
\end{enumerate}

Keep all decisions auditable. Write:

\begin{lstlisting}
windows.csv
clusters.csv
segments.csv
patches.csv
generation_summary.csv
\end{lstlisting}

Each generation should have a directory:

\begin{lstlisting}
generations/g00/
generations/g01/
...
\end{lstlisting}

\begin{center}\rule{0.5\linewidth}{0.5pt}\end{center}

\section{Rescue phase}\label{rescue-phase}

After the two ends meet/cross:

\begin{enumerate}
\def\labelenumi{\arabic{enumi}.}
\tightlist
\item
  Fit the global PMF using all current patches.
\item
  Find: \[
  x^* = \arg\max_x \mathrm{var}_{global}(x)
  \] among grid points with finite global PMF.
\item
  Add a rescue window near \(x^*\):

  \begin{itemize}
  \tightlist
  \item
    center at nearest grid point to \(x^*\)
  \item
    \passthrough{\lstinline!k = k\_rescue!} by default
  \item
    optionally increase \passthrough{\lstinline!k!} if the target is not
    sampled in the tail quantile band.
  \end{itemize}
\item
  Run EQ for the rescue window.
\item
  Recompute clusters/segments/patches.
\item
  Refit global PMF.
\item
  Repeat until:

  \begin{itemize}
  \tightlist
  \item
    max rescue rounds reached,
  \item
    total budget exhausted,
  \item
    or maximum global variance falls below a practical threshold if
    implemented.
  \end{itemize}
\end{enumerate}

Write:

\begin{lstlisting}
rescue_summary.csv
rescue_round_<n>/...
\end{lstlisting}

\begin{center}\rule{0.5\linewidth}{0.5pt}\end{center}

\section{Important behavior}\label{important-behavior}

\subsection{Do not hard-stitch PMFs}\label{do-not-hard-stitch-pmfs}

Do not implement final PMF as:

\begin{lstlisting}[language=Python]
if var_f[x] < var_r[x]:
    use forward
else:
    use reverse
\end{lstlisting}

Do not implement final PMF as choosing one patch per bin.

Use variance-weighted global fusion.

\subsection{Use minimum variance only to define patch
uncertainty}\label{use-minimum-variance-only-to-define-patch-uncertainty}

Allowed:

\begin{lstlisting}[language=Python]
var_EQ(x) = min_m var_m(x)
var_NEQ(x) = min(var_left_anchor(x), var_right_anchor(x))
\end{lstlisting}

Not allowed as the final PMF rule:

\begin{lstlisting}[language=Python]
PMF(x) = PMF_from_lowest_variance_patch(x)
\end{lstlisting}

\subsection{Keep local gauges separate from global
gauge}\label{keep-local-gauges-separate-from-global-gauge}

Each patch can have arbitrary local gauge. The global fusion step must
solve additive offsets.

\subsection{Write enough intermediate
files}\label{write-enough-intermediate-files}

The implementation should be easy to debug from disk outputs. For every
generation, write:

\begin{lstlisting}
window summaries
cluster summaries
segment summaries
patch PMF/variance files
global PMF file
\end{lstlisting}

\begin{center}\rule{0.5\linewidth}{0.5pt}\end{center}

\section{\texorpdfstring{2. Bash script:
\texttt{scripts/run\_mines\_variance\_fusion.sh}}{2. Bash script: scripts/run\_mines\_variance\_fusion.sh}}\label{bash-script-scriptsrun_mines_variance_fusion.sh}

Create a bash script that runs the new workflow with default parameters.

Usage:

\begin{lstlisting}[language=bash]
bash scripts/run_mines_variance_fusion.sh
\end{lstlisting}

Optional arguments:

\begin{lstlisting}[language=bash]
bash scripts/run_mines_variance_fusion.sh \
  --system-root results/doublewell_1d \
  --bin ./build/neq_sim \
  --seed 10101 \
  --label mines_variance_fusion_test
\end{lstlisting}

Defaults:

\begin{lstlisting}[language=bash]
SYSTEM_ROOT="results/doublewell_1d"
BIN="./build/neq_sim"
SEED=10101
LABEL="mines_variance_fusion"
TOTAL_BUDGET_STEPS=200000
T_NEQ=5000
N_NEQ_TRAJ=100
N_EQ_STEPS=10000
\end{lstlisting}

The script should:

\begin{enumerate}
\def\labelenumi{\arabic{enumi}.}
\tightlist
\item
  Check that the Python file exists.
\item
  Check that \passthrough{\lstinline!SYSTEM\_ROOT/run\_context.json!}
  exists.
\item
  Check that the binary exists and is executable.
\item
  Create log directory:
\end{enumerate}

\begin{lstlisting}
<SYSTEM_ROOT>/MINES/<LABEL>/logs/
\end{lstlisting}

\begin{enumerate}
\def\labelenumi{\arabic{enumi}.}
\setcounter{enumi}{4}
\tightlist
\item
  Run:
\end{enumerate}

\begin{lstlisting}[language=bash]
python scripts/mines_variance_fusion.py ...
\end{lstlisting}

\begin{enumerate}
\def\labelenumi{\arabic{enumi}.}
\setcounter{enumi}{5}
\tightlist
\item
  Tee stdout/stderr to a timestamped log file.
\end{enumerate}

\begin{center}\rule{0.5\linewidth}{0.5pt}\end{center}

\section{\texorpdfstring{3. Jupyter notebook:
\texttt{notebooks/mines\_variance\_fusion\_visualization.ipynb}}{3. Jupyter notebook: notebooks/mines\_variance\_fusion\_visualization.ipynb}}\label{jupyter-notebook-notebooksmines_variance_fusion_visualization.ipynb}

Create a notebook that visualizes the result.

It should allow the user to set:

\begin{lstlisting}[language=Python]
system_root = "../results/doublewell_1d"
label = "mines_variance_fusion"
seed = 10101
\end{lstlisting}

The notebook should load:

\begin{lstlisting}
<system_root>/MINES/<label>/raw/seed_<seed>/global_pmf.csv
windows.csv
clusters.csv
segments.csv
patches.csv
generation_summary.csv
rescue_summary.csv
\end{lstlisting}

Plots:

\subsection{Plot 1: Global PMF}\label{plot-1-global-pmf}

Plot:

\begin{lstlisting}
x vs global_pmf
x vs analytic_pmf
uncertainty band using ± sqrt(global_variance)
\end{lstlisting}

\subsection{Plot 2: Global variance}\label{plot-2-global-variance}

Plot:

\begin{lstlisting}
x vs global_variance
\end{lstlisting}

Mark rescue target positions if available.

\subsection{Plot 3: Window centers}\label{plot-3-window-centers}

Plot window centers over generation:

\begin{lstlisting}
generation vs center_x
\end{lstlisting}

Color or marker by side:

\begin{lstlisting}
left, right, rescue, bridge
\end{lstlisting}

\subsection{Plot 4: Patch coverage}\label{plot-4-patch-coverage}

For each patch, show which grid points it covers.

A simple heatmap is fine:

\begin{lstlisting}
x-axis: x
y-axis: patch name
value: coverage_mask
\end{lstlisting}

\subsection{Plot 5: Local patch PMFs}\label{plot-5-local-patch-pmfs}

Overlay all local patch PMFs after applying fitted offsets from
\passthrough{\lstinline!global\_fit\_summary.json!}.

\subsection{Plot 6: Error profile}\label{plot-6-error-profile}

If \passthrough{\lstinline!analytic\_pmf!} is available:

\begin{lstlisting}
x vs global_pmf - analytic_pmf
\end{lstlisting}

Also print RMSE over finite points.

The notebook should be robust if some files do not exist yet.

\begin{center}\rule{0.5\linewidth}{0.5pt}\end{center}

\section{Acceptance criteria}\label{acceptance-criteria}

The implementation is considered correct if:

\begin{enumerate}
\def\labelenumi{\arabic{enumi}.}
\tightlist
\item
  It creates
  \passthrough{\lstinline!scripts/mines\_variance\_fusion.py!}.
\item
  It creates
  \passthrough{\lstinline!scripts/run\_mines\_variance\_fusion.sh!}.
\item
  It creates
  \passthrough{\lstinline!notebooks/mines\_variance\_fusion\_visualization.ipynb!}.
\item
  The Python script only implements MiNES, not AUS.
\item
  EQ-overlapping ensembles are merged into EQ clusters.
\item
  EQ clusters use MBAR PMF reconstruction.
\item
  EQ cluster variance is computed as:
\end{enumerate}

\[
\mathrm{var}_{EQ}(x)=\min_m \mathrm{var}_m(x)
\]

where each \(\mathrm{var}_m(x)\) comes from bootstrap PMFs aligned to
\(x^{most}_m\).

\begin{enumerate}
\def\labelenumi{\arabic{enumi}.}
\setcounter{enumi}{7}
\tightlist
\item
  Non-overlapping neighboring ensembles/clusters use bidirectional
  NEQ/MTS.
\item
  NEQ variance is computed as:
\end{enumerate}

\[
\mathrm{var}_{NEQ}(x)=\min(\mathrm{var}_{left-anchor}(x), \mathrm{var}_{right-anchor}(x)).
\]

\begin{enumerate}
\def\labelenumi{\arabic{enumi}.}
\setcounter{enumi}{9}
\tightlist
\item
  The final PMF is not hard-stitched. It is obtained by inverse-variance
  weighted global fusion with fitted patch offsets.
\item
  The rescue target is selected from the maximum global fused variance.
\item
  All intermediate PMF, variance, window, cluster, segment, and rescue
  summaries are written to disk.
\item
  The bash script runs the workflow with sensible defaults.
\item
  The notebook visualizes global PMF, variance, windows, patches, rescue
  targets, and error against analytic PMF.
\end{enumerate}

\begin{center}\rule{0.5\linewidth}{0.5pt}\end{center}

\section{Notes}\label{notes}

Prioritize correctness and clarity over cleverness. Keep the file
modular. Avoid mixing old experimental protocols into this new
implementation. This should become the clean reference implementation of
MiNES with variance-weighted PMF fusion.

\end{document}
